\section{Motivation: Two Structural Limits}
\label{sec:motivation}

Packet capture has been shaped by tools that assume the kernel network
stack is a stable observation surface. That assumption no longer
holds. BPF programs deployed at XDP, TC, socket-filter and cgroup hooks
now intercept, drop, redirect, or rewrite packets long before they
reach the surfaces tools like tcpdump can see. We argue that
the resulting observability gap is structural, not a tooling defect,
and that closing it requires addressing two limits simultaneously.

\subsection{Two Structural Limits}
\label{sec:two-limits}

\paragraph{Limit~A: filter expressiveness.}
pcap-filter, the de-facto packet-filter language since 1993,
compiles to classic BPF (cBPF). cBPF has no loop primitive, by design:
the predecessor virtual machine guaranteed termination by forbidding
backward branches. That decision excludes the structures that
characterise modern protocols: SRv6 segment lists, IPv6 extension-header
chains, TCP-option walks, and multi-level tunnels (VXLAN, Geneve, GTP-U,
MPLS). Recent alternatives are partial fixes:
\texttt{tc-flower}~\cite{tcflower} matches a single encapsulation level
via \texttt{enc\_*} keys but offers no expression syntax;
\texttt{nftables}~\cite{nftables_wiki} is rule-based with rich L3/L4
plus connection-tracking but no deep encapsulation;
\texttt{NetPFL}~\cite{ciminiera2010netpfl} addressed tunnel-aware
filtering but targeted a userspace VM and has been inactive since
roughly 2013.

\paragraph{Limit~B: attachment-point coverage.}
The Linux network stack is now a programmable data path: BPF can
attach at XDP, TC ingress/egress, socket filter, cgroup/skb,
sock\_ops, sk\_lookup, netfilter BPF, flow\_dissector, and several
tracing hooks. Production deployments combine multiple hooks --
Cilium uses XDP, TC, sockops, and cgroup/skb together for Kubernetes
networking; Cloudflare's L4 load balancer~\cite{cloudflare_unimog} runs
on XDP; Meta's Katran~\cite{shukla2020katran} runs on XDP. Packets
that these programs \texttt{XDP\_DROP} never reach AF\_PACKET, so
tcpdump cannot see them. Existing XDP-aware tools either
lack filtering (\texttt{xdpdump}~\cite{xdptools_xdpdump}) or require
in-place hooks compiled into the production program
(\texttt{xdpcap}~\cite{cloudflare2019xdpcap}). \texttt{pwru}~\cite{cilium_pwru}
follows packets across kernel functions but produces a trace log, not
a pcap, and uses pcap-filter so does not address Limit~A.

Limits A and B are not independent: both stem from observability
infrastructure that pre-dates the BPF-programmable data path. Solving
either alone leaves a tool that knows where to attach but cannot
filter, or filters expressively but only at one hook.

\subsection{No Prior System Addresses Both}
\label{sec:no-prior}

\begin{table*}[t]
\centering
\caption{Capability matrix across packet-observation tools (\yes\ supported; \partsym\ partial; \no\ unsupported; \na\ N/A). No existing tool covers both structural limits.}
\label{tab:tools}
\small
\begin{tabular}{@{}l c c c c c c c@{}}
\toprule
                                & tcpdump & xdpcap & xdpdump & nft+nflog & tc-flower & pwru & \kunai+\xdpninja \\
\midrule
Output: pcap                    & \yes & \yes & \yes & \partsym & \no & \no & \yes \\
Visible to XDP\_DROP            & \no  & \partsym & \yes & \no & \no & \yes & \yes \\
Multi-encapsulation chain       & \no  & \no  & \na & \partsym & \partsym & \no & \yes \\
Variable-length lists / walks   & \no  & \no  & \na & \no & \no & \no & \yes \\
Expression syntax (where-clause)& \partsym & \partsym & \na & \partsym & \no & \partsym & \yes \\
Verifier-safe by construction   & \na  & \partsym & \na & \na & \na & \na & \yes \\
Non-invasive deployment         & \yes & \no  & \yes & \yes & \na & \yes & \yes \\
\bottomrule
\end{tabular}
\end{table*}

Table~\ref{tab:tools} compares seven tools across the two limits. No
single row covers both. The closest contenders are \texttt{xdpcap}
(addresses Limit~B for XDP but is invasive and uses cBPF) and
\texttt{tc-flower} (rich at TC but match-action, not capture, and
limited to single-level encapsulation). \texttt{pwru} sits in a
different category --- a kernel-wide tracer, complementary rather than
competing with capture (\S\ref{sec:related}).

\subsection{Variable-Length Headers Are Now the Norm}
\label{sec:varlen-norm}

To make Limit~A concrete, we sketch three filtering tasks that defeat
pcap-filter and modern alternatives but are one line in
\kunai. The kunai output is a verifier-safe eBPF program (\S\ref{sec:codegen}).

\paragraph{(a) SRv6 segment-list lookup.}
A 5G or SDN backbone operator wants packets whose Segment Routing
Header lists a specific SID, say \texttt{fc00::abcd}. pcap-filter
can identify routing-header type 4 but cannot search the segment array:
the list length is in the header and cBPF cannot loop. \kunai writes:
\begin{lstlisting}[style=kunai]
eth/ipv6/srv6 where any(srv6.segments.addr == fc00::1)
\end{lstlisting}

\paragraph{(b) GTP-U with optional fields.}
A 5G core operator wants inner IPv4 destination 10.0.0.1. The GTP
header carries optional fields gated by presence flags, so a
fixed-offset filter sees the wrong bytes. \kunai walks the gating:
\begin{lstlisting}[style=kunai]
eth/ipv4@outer/udp/gtp/ipv4@inner/tcp where inner.dst == 10.0.0.1
\end{lstlisting}

\paragraph{(c) Geneve overlay inner IPv4.}
Container networking (Cilium, GCP, Azure VNet) tunnels traffic in
Geneve. \texttt{tc-flower}'s \texttt{enc\_*} can match outer
encapsulation keys, but cannot follow Geneve's transparent inner
Ethernet to filter on inner L3:
\begin{lstlisting}[style=kunai]
eth/ipv4@outer/udp/geneve/eth/ipv4@inner/tcp where inner.dst == 10.0.0.1
\end{lstlisting}

These are not edge cases. SRv6 is deployed in 5G transport (RFC 8754);
Geneve and GTP-U are in production at every cloud and mobile
operator. The filtering need is decade-old; what changed is that the
toolchain has not kept pace.

\subsection{Roadmap}
\label{sec:roadmap}

The remainder of the paper addresses both limits as one architecture.
\S\ref{sec:xdp-ninja} introduces \xdpninja's non-invasive XDP attachment,
which closes Limit~B for XDP without modifying the production program;
the same fentry/fexit mechanism applies to TC, which we evaluate in
\S\ref{sec:eval-verifier}. \S\ref{sec:kunai} presents \kunai's
layered DSL with a static type system and a P4-derived protocol
vocabulary, addressing Limit~A. \S\ref{sec:codegen} establishes that
the generated eBPF is verifier-safe across Linux 6.1--6.18.
\S\ref{sec:eval} reports expressiveness, codegen overhead,
end-to-end throughput, and the cross-kernel verifier matrix.
Throughout, production BPF programs remain bit-identical: this is a
hard design constraint, not an aspiration, and we discuss its
trade-offs in \S\ref{sec:transport-tradeoff} and \S\ref{sec:discussion}.

\TODO{Figure 2 (XDP\_DROP illustrative): production XDP program
DROPs SYN-flood traffic; tcpdump sees nothing on AF\_PACKET
while \xdpninja captures the SYNs at the XDP entry. Useful early in
\S\ref{sec:two-limits} as a visual hook for Limit~B.}
